\begin{proposition}[The normal jets along cusp, tacnode and triple curves]\hypertarget{tag-01FJ}{}\label{cdp:st:C2d}
\StatementRecord{cdp:st:C2d}{proved}{01FJ}{\SeriesRef{cdp:eq:C.J.2}, \SeriesRef{cdp:eq:C.J.3}, \SeriesRef{cdp:st:A.2.4}}

Let $Y$ be a smooth complex threefold with $b_2(Y)=0$, let $S$ be a reduced Cartier divisor, and let $D\subset\operatorname{Sing}S$ be a smooth compact connected curve. Suppose the general plane section transverse to $D$ has one of the equations
\begin{equation}
 x^2-y^3=0,\qquad x^2-y^4=0,\qquad xy(x-y)=0.
 \hypertarget{tag-00OV}{}\label{cdp:eq:C.J.1}
\end{equation}
Then $g(D)\le1$. For $g(D)=1$, the indicated transverse analytic type is constant along all of $D$. The inverse image of $D$ in the normalization is respectively an isomorphic copy, a finite \'etale cover of degree two, or a finite \'etale cover of degree three. Each connected component $C$ of that cover has genus one and normal line of degree zero in the normalization, which is smooth near $C$. At these three transverse types, the inverse Jacobian ideal pulls back into the line allowing at most a simple pole on each normalization branch.
\end{proposition}


\begin{proof}

\begin{enumerate}[label={$\langle1\rangle$\arabic*.},leftmargin=*]
\item Put $N=N_{D/Y}$ and $B=\OO_Y(S)|_D$. Since every ambient line bundle has zero rational first Chern class, $\deg B=\deg K_Y|_D=0$. Adjunction for the smooth curve gives
\begin{equation}
 \det N=K_D\otimes K_Y^{-1}|_D,\qquad
 \deg N=2g(D)-2.                                      \hypertarget{tag-00OW}{}\label{cdp:eq:C.J.2}
\end{equation}
A local defining section of $S$ belongs to $I_D^2$. Its quadratic normal jet is a well-defined section of $\operatorname{Sym}^2N^*\otimes B$. In the first two cases its generic rank is one. Its saturated radical $R\subset N$ is a line subbundle on the smooth curve, and $Q=N/R$ is a line bundle. The quadratic jet gives a nonzero section of $B\otimes Q^{-2}$, so
$\deg Q\le0$.

\item In the cusp case, the cubic jet on the radical is an intrinsic nonzero section of $B\otimes R^{-3}$. To check its independence, choose coordinates $(x,y,z)$ with $D=(x=y=0)$ and $R$ the $y$ direction. A change of lift of a radical vector changes the cubic by a Hessian term with one radical argument, which is zero. Changing the local retraction to $D$ differentiates the identity that the quadratic jet vanishes on $R$, again giving zero. Thus the local cubics glue across all points, including possible zeros of the quadratic coefficient. Consequently $\deg R\le0$. Equation \eqref{cdp:eq:C.J.2} proves $g(D)\le1$.

\item In the tacnode case the radical cubic vanishes identically. Write the relevant normal terms as
\[
 a x^2+bxy^2+cy^4.
\]
Completing the square has invariant numerator
\begin{equation}
 \Delta=ac-b^2/4\in
 H^0(D,B^2\otimes Q^{-2}\otimes R^{-4}).               \hypertarget{tag-00OX}{}\label{cdp:eq:C.J.3}
\end{equation}
This numerator is holomorphic even at a zero of $a$. Under $x\mapsto x+uy^2$, the coefficients become $b+2au$ and $c+bu+au^2$, and \eqref{cdp:eq:C.J.3} is unchanged. Linear changes preserving the radical give precisely the displayed line-bundle weights. Changes of the retraction to $D$ add terms containing the radical cubic or a radical Hessian argument, so also leave this quartic numerator unchanged. These checks give its global definition. It is nonzero by the generic tacnode assumption. Hence
\[
 \deg Q\le0,\qquad \deg Q+2\deg R\le0,
 \qquad \deg N\le\tfrac12\deg Q\le0.
\]
This proves the same genus bound.

\item For the ordinary triple case the quadratic jet is identically zero. The binary cubic normal jet is a section of $\operatorname{Sym}^3N^*\otimes B$. Its discriminant is a nonzero section of
$B^4\otimes(\det N^*)^6$. The discriminant has degree four in the cubic coefficients and weight six under a linear change of the two normal coordinates; these weights also follow by writing it as the square of the product of the three pairwise determinants of the linear factors. Its nonnegative divisor degree gives $\deg N\le0$ and the genus bound again.

\item When $g(D)=1$, all the inequalities are equalities. In the cusp case both the quadratic and radical cubic sections have no zeros. In the tacnode case $a$ and $\Delta$ have no zeros. In the triple case the cubic discriminant has no zeros. The holomorphic splitting lemma with parameter $z$ completes the square in the first two cases. The remaining one-variable term is respectively a unit times $y^3$ or $y^4$; taking a local root of the unit gives the corresponding equation in \eqref{cdp:eq:C.J.1}. In the triple case its three distinct tangent lines split after an \'etale ordering cover; the holomorphic implicit function theorem gives three smooth branches meeting transversely in each normal plane. Straightening two branches to the axes and rescaling the third branch's unit gives the third equation. Thus these are constant transverse types throughout $D$. In particular no other singular curve meets $D$.

\item In the cusp normalization, $x=u^3,y=u^2$, up to the choice of which coordinate is the radical. The leading terms $y=u^2$ and $x=u^3$ identify
$P^2$ with $R$ and $P^3$ with $Q$, where
$P=N_{C/\widetilde S}$. Adapted coordinate changes preserve these
leading terms through the corresponding line transitions. Thus
$\deg P=0$. For the tacnode, the two branches have $x=\pm y^2$ and normal line the pullback of $R$, again of degree zero. Their permutation cover is \'etale because $\Delta$ is nowhere zero. For the ordinary triple, order the three tangent lines by a finite \'etale cover. Three pairwise transverse line subbundles in a rank-two bundle make it a sum $P\oplus P$: the third is a graph of an isomorphism between the first two. Since the determinant has degree zero, $\deg P=0$. Each normalization branch has this tangent-root normal line. \'{E}tale descent gives the claimed degrees and covers before ordering.

\item Compute the inverse Jacobian ideal in the parameter-preserving
normal forms, with coefficients holomorphic in the curve parameter $z$.
A biholomorphic coordinate change transforms the gradient by its
invertible Jacobian matrix; multiplication of the defining equation
by a unit multiplies the gradient by that unit modulo the equation.
These maps transport the following fractional ideals to $S$.

For the cusp set $B=\C\{u,z\}$ and
$A=\C\{u^2,u^3,z\}=\C\{z\}+u^2B$.
The Jacobian ideal has images $u^3,u^4$. For a multiplier $h$, put
$a=u^3h$. The conditions $a,ua\in A$ kill respectively the linear
and constant coefficients of $a$. Hence
$\mathcal J^{-1}=u^{-1}B$, agreeing with
Proposition~\ref{cdp:st:A.2.4} under the displayed normalization map.

For the tacnode the two maps are $x=\pm u^2,y=u$, and the image ring is
\[
 A=\{(f,g)\in B^2:f-g\in u^2B\}.
\]
Indeed, the inverse expression is
$(f+g)/2+x(f-g)/(2u^2)$. The Jacobian generators have images
$(u^2,-u^2)$ and $(u^3,u^3)$, up to nonzero scalar factors.
For $a=u^2h_1$ and $b=u^2h_2$ the membership conditions are
$a+b\in u^2B$ and $a-b\in uB$. Their constant coefficients vanish,
so each $h_i$ has at most a simple pole. Their residues are opposite:
\[
 \mathcal J^{-1}=B^2+\C\{z\}(u^{-1},-u^{-1}).
\]

For the ordinary triple point use the branches
$(0,u),(u,0),(u,u)$. The image ring consists of triples with a common
constant and with third linear coefficient equal to the sum of the
first two. All coefficients of degree at least two are independent:
for $k\ge2$ the polynomials
$y^k-xy^{k-1}$, $x^k-x^{k-1}y$, and $x^{k-1}y$ have respectively
the three individual order-$k$ branch values.
The Jacobian generators have images
$(-u^2,0,u^2)$ and $(0,u^2,-u^2)$.
For $a_i=u^2h_i$, the common-constant condition on both products forces
all three constant coefficients to vanish. The linear relations
give residues $(a,a,-a)$. Thus
\[
 \mathcal J^{-1}=B^3+\C\{z\}(u^{-1},u^{-1},-u^{-1}).
\]
For the tacnode and triple point, regularity of the products on each
normalization branch excludes poles of order greater than two before
these coefficient conditions are imposed. The three displayed modules are contained in the
invertible sheaf allowing a simple pole on the reduced inverse-image
curve. This proves the claimed bound at every point.

\end{enumerate}
\end{proof}